\documentclass[11pt,a4paper]{article}
\usepackage[utf8]{inputenc}
\usepackage[T1]{fontenc}
\usepackage[margin=2cm]{geometry}
\usepackage{amsmath,amssymb}
\setlength{\parindent}{0pt}
\setlength{\parskip}{6pt}

\begin{document}

\begin{center}
{\LARGE\textbf{Cahier de Physique I}}\\
\large Applications physiques du crible congruentiel : vide adélique,\\
renormalisation dyadique, chaos quantique, matière condensée\\
Fouad Bensmail --- 25 septembre 2026
\end{center}

Ce cahier est le frère physique du cahier mathématique \emph{Addendum v4 à la
Note Universelle VII} (47 addenda, Zenodo v39, DOI 10.5281/zenodo.22950183).
Il en applique les mesures --- volumes locaux, nombres de Tamagawa, résidus
spectraux, grammaire dyadique --- aux quatre portes physiques ordonnées par
degré de résonance : mécanique quantique adélique et cordes $p$-adiques
(rang I) ; groupe de renormalisation, verres de spin et réseaux de tenseurs
(rang II) ; chaos quantique et systèmes ouverts (rang III) ; matière
condensée et transitions de phase (rang IV). Même méthode que la maison :
toute affirmation est une mesure avec garde ; tout refus est publié.

\vspace{0.5cm}
\hrule
\vspace{0.5cm}

\section*{Addendum I-Phys --- adelic\_amplitude : le vide adélique, point fixe infrarouge et séparation géométrie/matière}

Le compagnon \texttt{adelic\_amplitude.py} (premier du Cahier Physique) teste
l'hypothèse de Freund--Witten (1987) : les amplitudes de cordes $p$-adiques
s'adélisent en une constante universelle. Notre ``joyau'' (Addendum XIII :
$\prod_p \beta_p(\text{jumeaux}) = 2C_2$) est-il cette constante de
couplage ?

\textbf{Refus publié (v0).} Le protocole initial construisait l'amplitude
$p$-adique naïve $A_p(s) = \beta_p(H)\,p^{-s}$ et cherchait ses pôles :
objet mal défini --- le facteur $\prod_p p^{-s} = e^{-s\,\theta(P)}$
s'effondre vers $0$ pour tout $\mathrm{Re}\,s > 0$ ; renormalisé de cette
divergence triviale de volume, toute dépendance en $s$ disparaît. La v1 ne
mesure donc que trois objets exacts : le produit eulérien (le vide), sa
queue (le flot RG), son résidu spectral (la séparation géométrie/matière).

\textbf{Protocole.} Pour six motifs admissibles $H$, on calcule les volumes
locaux $\beta_p(H)$ pour $p \le 10^7$, le produit partiel $S_P(H)$, et la
queue eulérienne $S(H) - S_P(H) \approx C_H/(P \ln P)$. Trois objets sont
mesurés : (1) le point fixe infrarouge $S(H)$ ; (2) le flot de
renormalisation de la couplance adélique ; (3) le spectre FFT du résidu
eulérien sur $t \in [10, 50]$, en variable $\ln P$.

\textbf{Gardes pré-enregistrées.}
\begin{itemize}
\item G0 (couplage) : $|S_{10^6}(\text{jumeaux}) - 1{,}3203236| < 10^{-6}$ --- tenue.
\item G1 (flot RG) : variation relative de $C_H(P) = (S_{\text{ancre}} - S_P)\,P\ln P$ sur $[10^4, 10^6] < 15\,\%$ pour chacun des six motifs --- tenue ($-7{,}0\,\%$ partout).
\item G3 (fonction bêta) : pente de $\ln|S_{\text{ancre}} - S_P| + \ln\ln P$ contre $\ln P$ dans $[-1{,}15 ; -0{,}85]$ --- tenue ($-1{,}007$ partout).
\item G2 (séparation) : aucun pic FFT du résidu eulérien $> 3\times$ la médiane sur $t \in [10, 50]$ $\to$ vide spectralement silencieux (prédit) ; un pic à $< 0{,}45$ d'un zéro de $\zeta$ $\to$ STRUCTURE (ouverture) --- tenue ($1{,}00$ partout).
\end{itemize}

\textbf{Mesure.} Temps total : $1{,}47$ s.

\begin{center}
\begin{tabular}{lcccc}
\hline
Motif & $S_{10^6}$ & $C_H$ & var (\%) & pente \\
\hline
jumeaux    & 1{,}3203237 & $-1{,}2046$  & $-7{,}0$ & $-1{,}007$ \\
cousins    & 1{,}3203237 & $-1{,}2046$  & $-7{,}0$ & $-1{,}007$ \\
sexy       & 2{,}6406474 & $-2{,}4092$  & $-7{,}0$ & $-1{,}007$ \\
triplet\_A & 2{,}8582492 & $-7{,}8234$  & $-7{,}0$ & $-1{,}007$ \\
triplet\_B & 2{,}8582492 & $-7{,}8234$  & $-7{,}0$ & $-1{,}007$ \\
quadruplet & 4{,}1511826 & $-22{,}7249$ & $-7{,}0$ & $-1{,}007$ \\
\hline
\end{tabular}
\end{center}

Spectre FFT du résidu eulérien (jumeaux) aux neuf zéros de $\zeta$ :
amplitude $0{,}314$ à $t = 14{,}13$ et $0{,}637$ aux huit autres --- toutes
au plancher de bruit ($\max/\text{médiane} = 1{,}00$ sur les six motifs).

\textbf{Verdict.} \texttt{AUDIT VERT}. Gardes G0, G1, G2, G3 tenues.

\textbf{Lecture.} La séparation géométrie/matière est mesurée :
\begin{enumerate}
\item Le vide eulérien (produit pur des $\beta_p$) est un point fixe
stationnaire, spectralement silencieux : il ne porte pas les résonances de
$\zeta$ détectées dans le résidu du crible (Addenda XXXVI--XLVII). En langue
physique : le fond géométrique ne porte pas les fluctuations quantiques du
spectre arithmétique.
\item Le flot RG de la queue eulérienne est exactement en $1/(P \ln P)$
(pente $-1{,}007$), comme prédit par le Théorème T3 : la fonction bêta du
vide adélique est mesurée.
\item Le nombre de Tamagawa $2C_2$ est la constante de couplage infrarouge
du vide, confirmée à $10^{-7}$ près pour les jumeaux : le ``joyau'' de
l'Addendum XIII reçoit sa lecture physique de point fixe.
\end{enumerate}

\textbf{Statut.} Porte 2 (mécanique quantique adélique) ouverte par une
première pierre verte : le vide adélique est spectralement silencieux, et
les résonances $\zeta$ vivent dans le secteur matière (résidu du crible),
non dans le secteur géométrie (queue eulérienne). La frontière entre le
fond géométrique et les excitations qui y vivent est mesurée, non postulée.

\vspace{1cm}
\hrule
\vspace{1cm}

\section*{Addendum II-Phys --- dyadic\_renormalization : flot de renormalisation 2-adique et point fixe infra-rouge}

Le compagnon \texttt{dyadic\_renormalization.py} teste l'analogie entre la grammaire dyadique des premiers (Addenda XLII, XLIV) et le flot de renormalisation (RG) dans les systèmes hiérarchiques (verres de spin, réseaux MERA). La profondeur $k = v_2(p-1)$ est l'échelle de RG ; la partie impaire $m = (p-1)/2^k$ représente les degrés de liberté résiduels.

\textbf{Protocole.} Pour $p \le 10^7$, on regroupe les $m$ par profondeur $k$. Pour chaque $k$, on calcule un \textit{score de régularité} (inverse de la variance normalisée des écarts entre $m$ triés), comparé à un null stratifié (échantillons aléatoires de même taille).

\textbf{Gardes pré-enregistrées.}
\begin{itemize}
\item G1 (Enchevêtrement) : Pour $1 \le k \le 8$, le score doit être $> 3\sigma$ au-dessus du null.
\item G2 (Saturation / Point fixe) : Pour $k \ge 10$, le score doit cesser de croître (stagnation ou effondrement).
\item G3 (Robustesse d'échelle) : La forme de la courbe de régularité doit être fortement corrélée (Spearman $\rho > 0.8$) entre la fenêtre $P \le 10^6$ et $P \le 10^7$.
\end{itemize}

\textbf{Mesure.} Temps d'exécution : $0.79$ s.
\begin{itemize}
\item \textbf{G1 : TENUE.} La régularité est massivement significative (ex: $+154.1\sigma$ à $k=1$, $+13.4\sigma$ à $k=8$).
\item \textbf{G2 : TENUE.} Le score atteint un pic à $k=11$ ($1.59$) puis s'effondre à $k=13$ ($1.40$), validant la saturation hiérarchique.
\item \textbf{G3 : REFUS PUBLIÉ.} Corrélation de Spearman $\rho = 0.297$.
\end{itemize}

\textbf{Analyse du refus (G3).} L'échec de la robustesse d'échelle n'invalide pas G1 et G2. Il s'agit d'un \textit{artefact de taille finie} : à $P = 10^6$, les bins de grande profondeur ($k \ge 10$) contiennent trop peu d'éléments ($< 50$) pour une estimation stable de la variance. Ce bruit statistique détruit la corrélation de forme avec la courbe lisse de $10^7$. 

\textbf{Lecture physique.} L'enchevêtrement dyadique (ordre croissant avec l'échelle) et la saturation (point fixe infra-rouge) sont massivement confirmés. La métrique de régularité est valide, mais sa robustesse nécessite des fenêtres d'observation au-delà du seuil de bruit statistique des petits échantillons.

\textbf{Statut.} Porte 2 (Rang II) consolidée. Le phénomène de renormalisation 2-adique est mesuré ; la limite de la métrique est cartographiée et publiée.

\vspace{1cm}
\hrule
\vspace{1cm}

\section*{Addendum III-Phys --- quantum\_scars (v1) : résonances transitoires et cicatrices quantiques}

Le compagnon \texttt{quantum\_scars.py} (v1) teste l'analogie entre le résidu spectral du crible et les systèmes quantiques ouverts. En chaos quantique, les systèmes ouverts ne possèdent pas d'états liés stables mais des \textit{états de Gamow} (résonances à durée de vie finie, ajustables par un profil de Lorentz). De même, les systèmes chaotiques présentent des \textit{cicatrices quantiques} (quantum scars), des renforcements locaux de probabilité le long d'orbites instables, sans former de réseau périodique rigide.

\textbf{Protocole (v1).} Sur une fenêtre glissante de premiers ($P \in [10^6, 2 \cdot 10^6]$, taille $W=50\,000$), on calcule la densité locale discrétisée (histogramme) et son spectre de puissance (FFT). On cherche des pics spectraux transitoires dépassant $3\sigma$ par rapport à un null stratifié généré par \textit{surrogats de phase} (randomisation des phases de la FFT).

\textbf{Gardes pré-enregistrées.}
\begin{itemize}
\item G1 (Non-stationnarité) : Le coefficient de variation (CV) de l'amplitude des pics à travers les fenêtres doit être $> 0.15$.
\item G2 (Profil de Résonance) : Le pic le plus significatif doit s'ajuster à un profil de Lorentz avec un $R^2 > 0.85$.
\item G3 (Absence de Cristal) : La distribution des fréquences des pics significatifs ne doit pas former un réseau périodique (CV des espacements $> 0.3$).
\end{itemize}

\textbf{Mesure.} Temps d'exécution : $0.06$ s.
\begin{itemize}
\item \textbf{G1 : TENUE.} CV $= 0.161$. L'amplitude spectrale fluctue, confirmant la non-stationnarité du signal.
\item \textbf{G2 et G3 : REFUS PUBLIÉ.} Aucun pic spectral ne dépasse le seuil de $3\sigma$ par rapport aux surrogats de phase. L'ajustement lorentzien et le test de périodicité sont donc impossibles.
\end{itemize}

\textbf{Analyse du refus (Diagnostic de l'observable).}  
L'échec de G2 et G3 n'invalide pas l'hypothèse physique, mais révèle une limite de l'observable choisie en v1. La densité locale des nombres premiers (discrétisée par histogramme sur des fenêtres de $5 \cdot 10^4$) est dominée par le bruit de Poisson intrinsèque à la distribution. Les surrogats de phase ont correctement identifié ce bruit de fond : les fluctuations spectrales observées ne sont pas des résonances structurelles, mais des artefacts de la discrétisation grossière. 

En physique des systèmes chaotiques et dans l'étude des zéros de $\zeta$, les signatures spectrales fines (résonances de Gamow, statistique de Wigner-Dyson) ne se lisent pas dans la densité spatiale brute, mais dans la \textbf{distribution des écarts entre niveaux consécutifs} (les \textit{gaps} $g_n = p_{n+1} - p_n$) ou via des analyses multi-échelles (ondelettes) sur des fenêtres beaucoup plus vastes.

\textbf{Lecture physique.} La non-stationnarité du résidu est confirmée (G1), mais l'observable ``densité locale sur fenêtre courte'' est un microscope insuffisant pour révéler la structure fine des résonances transitoires. Le bruit de fond arithmétique masque le signal spectral. Une version ultérieure (v2) devra cibler la distribution des écarts (gaps) ou utiliser une résolution adaptative pour franchir le plancher de bruit de Poisson.

\textbf{Statut.} Porte 3 (chaos quantique) entrouverte. La non-stationnarité est mesurée ; la limite de l'observable de densité est cartographiée et publiée.

\vspace{1cm}
\hrule
\vspace{1cm}

\section*{Note Méthodologique et Positionnement Bibliographique}

\textbf{1. Consolidation numérique par haute précision.}  
Lors de la mesure initiale de la queue eulérienne, une anomalie numérique a été détectée pour les motifs de taille $k \ge 3$ : l'exposant effectif $\alpha_{\text{eff}}$ de la décroissance de l'erreur semblait diverger à $P = 10^6$. Le diagnostic a révélé un artefact instrumental : les constantes d'ancre $\Theta_\infty$ étaient codées avec 7 décimales, créant un plancher de bruit. Le script \texttt{alpha\_convergence\_mp.py} a corrigé cela en utilisant la bibliothèque \texttt{mpmath} (50 décimales) et des ancres calculées à $P = 10^7$. L'anomalie a disparu, confirmant que $\alpha_{\text{eff}}$ converge uniformément vers $1 + 1/\ln P$, validant la décroissance en $1/(P \ln P)$ prédite par le Théorème T3.

\textbf{2. Recoupements bibliographiques et positionnement (Article III).}  
Ce travail ne prétend pas à la découverte de nouveaux théorèmes analytiques, mais à la construction d'un laboratoire numérique rigoureux qui vérifie, mesure et cartographie la frontière du prouvable. Les recoupements suivants confirment la validité de nos mesures :
\begin{itemize}
\item \textbf{Formule de la queue eulérienne $C(k) = \Theta_\infty \cdot k(k-1)/2$ :} Cette formule, mesurée empiriquement ici, est un cas particulier du théorème général de Bateman-Horn. Volfson (2026, \emph{Estimating the tail of the singular product}) établit ce coefficient via le groupe de Galois. Notre apport est la vérification numérique pour les constellations ($k=2,3,4$) et la mesure explicite de la correction de Mertens.
\item \textbf{Non-stationnarité spectrale et zéros de $\zeta$ :} Notre refus XLVII (les pics spectraux sont des résonances transitoires, pas un cristal) trouve un écho dans Goldston et Suriajaya (2021, \emph{Acta Arithmetica}), qui expriment le terme d'erreur de la moyenne de Riesz de la série singulière en fonction des zéros de $\zeta$. Notre mesure du résidu spectral en est la manifestation empirique.
\item \textbf{Théorème CRT exact (T2) :} Krause (2026) formule notre théorème T2 dans un cadre géométrique de polytope croisé. Son ``axiome d'invariance de constellation'' est une version géométrique de notre invariance par translation et homothétie (T1).
\item \textbf{Problème de parité :} Notre diagnostic confirmant que le crible de dimension 2 est bloqué, tandis que la dimension 3+ ouvre des portes partielles sans franchir le mur, est partagé par la communauté (Selberg, Tao, Maynard).
\end{itemize}

\textbf{3. Contribution spécifique de ce laboratoire.}  
La littérature fournit des théorèmes ; nous fournissons des vérifications expérimentales et des courbes de convergence. Notre apport réside dans : (a) des mesures numériques reproductibles à plusieurs échelles ($10^4$ à $10^8$) avec gardes pré-enregistrées ; (b) une méthode de laboratoire stricte (carte prédit, console confirme, écart publié) ; (c) la cartographie honnête de la frontière (4 théorèmes inconditionnels, 1 conjecture mesurée, et au-delà, le silence accepté).

\end{document}